% future/fp.tex
% mainfile: ../perfbook.tex
% SPDX-License-Identifier: CC-BY-SA-3.0

\section{用于并行化的函数式编程}
\label{sec:future:Functional Programming for Parallelism}
%
\epigraph{函数式编程在并行应用中的奇怪失败。}
	 {Malte Skarupke}

当我在1980年代初上第一门functional programming课时，
教授断言无副作用的functional programming风格非常适合
轻松实现并行化和分析。
三十年后，这一断言依然存在，但主流生产环境中
并行函数式语言的使用微乎其微——这种状况可能与教授的另一个断言
（即程序既不应该维护状态也不应该进行I/O）并非完全无关。
Erlang等函数式语言有小众应用，
多线程支持也已添加到其他几种函数式语言中，
但主流生产使用仍然是C、C++、Java和Fortran等过程式语言的领地
（通常与OpenMP、MPI或coarray配合使用）。

这种情况自然引出一个问题：``如果分析是目标，
为什么不在分析之前将过程式语言转换为函数式语言？''
当然，对这种方法有许多反对意见，其中我列出三条：

\begin{enumerate}
\item	过程式语言经常大量使用全局变量，
	这些变量可以被不同的函数独立更新，
	或者更糟糕的是，被多个线程独立更新。
	注意，Haskell的\emph{monad}就是为了处理
	单线程全局状态而发明的，而多线程访问
	全局状态对函数式模型造成了更多的伤害。
\item	多线程过程式语言经常使用锁、原子操作和事务等
	同步原语，这些对函数式模型造成了更多的伤害。
\item	过程式语言可以对函数参数进行\emph{别名化}，
	例如，通过将指向同一结构的指针作为两个不同的参数
	传递给同一函数的同一次调用。
	这可能导致函数在不知情的情况下通过两个不同的
	（可能重叠的）代码序列更新该结构，
	这极大地使分析变得复杂。
\end{enumerate}

当然，鉴于全局状态、同步原语和别名的重要性，
聪明的functional programming专家已经提出了许多尝试
来调和函数式编程（functional programming）模型与它们之间的矛盾，
monad只是其中一个例子。

另一种方法是将并行过程式程序编译为函数式程序，
然后使用functional programming工具来分析结果。
但完全可以比这做得更好，因为任何真实的计算
都是一个具有有限输入并运行有限时间的大型有限状态机。
这意味着任何真实程序都可以转换为一个表达式，
尽管可能是一个大得不切实际的表达式~\cite{VijayDSilva2012-sas}。

然而，并行算法的许多低级内核可以转换为
足够小的表达式，可以轻松装入现代计算机的内存中。
如果这样的表达式与一个断言相结合，检查该断言是否会触发
就变成了一个可满足性问题。
尽管可满足性问题是NP完全的，但通常可以在
比生成完整状态空间所需的时间短得多的时间内解决。
此外，求解时间似乎只弱依赖于底层内存模型，
因此运行在弱排序系统上的算法也可以被检查~\cite{JadeAlglave2013-cav}。

一般方法是将程序转换为静态单赋值（SSA）形式，
使得对变量的每次赋值都创建该变量的一个单独版本。
这适用于所有活跃线程的赋值，因此生成的表达式
包含了所讨论代码的所有可能执行。
添加一个断言就相当于询问输入和初始值的任何组合
是否可能导致断言触发，如上所述，
这正是可满足性问题。

一个可能的反对意见是它不能优雅地处理任意循环构造。
然而，在许多情况下，可以通过将循环展开有限次数来处理。
此外，也许某些循环也可以通过归纳方法来折叠。

另一个可能的反对意见是自旋锁涉及任意长的循环，
任何有限的展开都无法捕获自旋锁的完整行为。
事实证明这个反对意见很容易克服。
不是建模完整的自旋锁，而是建模一个trylock，它尝试
获取锁，如果不能立即获取就中止。
然后必须精心构造断言，以避免在自旋锁因锁不可立即获取
而中止的情况下触发。
因为逻辑表达式与时间无关，所有可能的并发行为
都将通过这种方法被捕获。

最后一个反对意见是，这种技术不太可能能够处理
像Linux内核这样数百万行代码构成的完整规模的软件工件。
这很可能是事实，但事实仍然是，对Linux内核中
每个小得多的并行原语进行穷举验证将是相当有价值的。
事实上，推动这一方法的研究人员已经将其应用于
非平凡的实际代码，包括Linux内核中的Tree RCU
实现~\cite{LihaoLiang2016VerifyTreeRCU,MichalisKokologiannakis2017NidhuggRCU}。

这种技术的适用范围有多广还有待观察，但它是
形式化验证领域中更有趣的创新之一。
虽然functional programming的倡导者终于可能在其
函数式编程必将统治一切的断言上是正确的，但显然
这种长期被吹捧的方法论正在其形式化验证的主场
开始面临可信的竞争。
因此，仍然有理由怀疑functional programming统治的不可避免性。
